\section{Introduction}
\label{sec:intro}

\textbf{The rank-recruitment question.} Descriptive work on linear-attention
and fast-weight LLMs shows that the state's effective rank is bounded by
sequence length and by how many independent associations the state must
hold \citep{nazari2026rank, sun2026staterank}. A companion causal result on
this same DeltaNet testbed \citep{schlag2021linear} showed something
stronger: when a task \emph{forces} the state to solve $K$ independent
associative bindings under a hard single-state bottleneck, gradient descent
recruits effective rank $\approx K$ and the resulting write operator is
exact, not merely approximately correct, for a wide range of $K$. That raises
the natural next question this paper answers: \textbf{what happens as $K$
approaches the state's own dimensionality $\dstate$?} A linear memory has a
hard information-theoretic ceiling near $K \approx \dstate$; but a
\emph{trained} memory need not approach that ceiling gracefully, and prior
work in this testbed lineage already hinted it does not
\citep{barnfield2026sharp}. This paper measures exactly how the trained
system fails as load increases, locates the failure precisely, tests whether
it is fundamental to the architecture or an artifact of one scale, and then
tests the single most natural mechanistic account of it directly.

\textbf{Contribution 1 --- the cliff, located.} At $d_{\mathrm{state}}=64$,
held-out compositional recovery ($h4 = \mathrm{rec@0.9}$ at hop depth 4) does
not decline smoothly with $K/\dstate$ --- it transitions sharply. We fit a
logistic to six points spanning $K/\dstate \in [0.25, 0.75]$ and locate the
transition's midpoint to $x_0 = 0.5455$ (95\% CI $[0.5385, 0.5513]$, width
$0.0127$, characteristic width $w=0.0597$), a $\sim\!20\times$ tightening
over the prior 0.25-wide bracket between the two flanking anchors.

\textbf{Contribution 2 --- the cliff, dissolved.} Re-running the identical
$K/\dstate$ window at $d_{\mathrm{state}}=128$ shows no cliff at all:
$h4=1.0$ at all four matched loads, all seeds, $12/12$ cells. The transition
is not shifted --- it exits the measured window entirely. This rules out a
naive ``$K/\dstate$ ratio alone determines the cliff'' account.

\textbf{Contribution 3 --- the leading confound, tested and exonerated.}
The two waves above differ in more than raw scale: the $d=64$ table (107
entities packed into 64 dimensions) is forced non-orthogonal by the Welch
bound, while the $d=128$ table (107 entities, 128 dimensions) is exactly
orthogonal. This is the single most natural remaining mechanistic
candidate --- and we test it directly, not by inference. Injecting
controlled, frozen coherence into the otherwise-flat $d=128$/$K=68$ cell,
swept from $0$ to $0.40$ (exceeding $d=64$'s own trained coherence band),
produces no cliff at any dose under a concentrated (rank-4) injection
structure: $h4=1.0$, $10/10$ cells, no exception.

\textbf{What this paper is NOT.} Not a claim that coherence-as-a-scalar can
never matter --- only that this specific, concentrated (rank-4) injection,
at this specific geometry, does not reproduce the cliff; a diffuse
injection structure is a registered, not-yet-run follow-on
(\S\ref{sec:open-question}). Not a resolution of \emph{why} the cliff
occurs at $d=64$ specifically --- this paper locates and characterizes the
transition and rules out one candidate mechanism; it does not close the
mechanism question. Not a claim about production-scale Gated-DeltaNet
\citep{siems2025deltaproduct, grazzi2025negative} or matrix-valued RNNs at
scale \citep{mishra2026m2rnn} --- all experiments here use a $\sim$sub-1M
parameter, single-layer synthetic-grammar testbed with 107 trainable
entities.
